\section{A local noisy many-neuron reversal certificate}
\label{sec:local-result}
The following result is a conditional entry theorem. Its hypothesis specifies
a neighborhood of an ID-trained state; it does not assume an OOD derivative
sign or a crossing. The local dynamics and the OOD conclusions are certified
below. Initialization-to-entry is a separate unresolved inequality.

Let $X$ be the supplied canonical binary dataset, $h=200$, and
$\widehat\theta_a$ the supplied numerical center indexed by $a=3.4$.
The loss uses all $N=4000$ samples. The fixed cohort $I$ is the original
$71$-neuron offline strong cohort. Let
\[
 \xi_\alpha=(\cos\alpha,\sin\alpha),\qquad
 \mathcal I=\{\alpha:|\alpha+17\pi/36|\le10^{-6}\}.
\]
The sector has angular width $2\cdot10^{-6}$ radians.

\begin{theorem}[Reversal from a certified learned-state neighborhood]
\label{thm:local-certificate}
Assume a compatible gradient-flow solution satisfies
\begin{equation}\label{eq:entry-missing}
 \|\theta(3.4)-\widehat\theta_{3.4}\|_2\le10^{-6}.
\end{equation}
Then, uniformly for $\alpha\in\mathcal I$,
\begin{align*}
 E(\xi_\alpha,3.4)-E(\xi_\alpha,3.7)&>1.2\cdot10^{-4},\\
 E(\xi_\alpha,3.9)-E(\xi_\alpha,3.7)&>4\cdot10^{-5}.
\end{align*}
In particular the full-network error attains an interior minimum on
$[3.4,3.9]$. At entry the fixed cohort has $m_{11}>0.9322$, input angles
below $15^\circ$, and output angles below $20^\circ$.
For the central probe, $Q_G<0$ at entry and $Q_G>0$ at $3.9$ for every
compatible training selector; the full-error rate has the same strict signs.
More quantitatively, for every probe in $\mathcal I$,
$\dot E<-0.00115$ almost everywhere on $[3.4,3.40001]$, and
$\dot E>0.000444$ almost everywhere on $[3.9,3.90001]$.
No continuous or unique zero of $Q_G$ is asserted.
\end{theorem}
\begin{proof}[Proof and computational obligations]
The local enclosure uses $2500$ cubic-reference steps of length $1/5000$.
At each step, outward bounds for the uniform reference defect, smooth
one-sided rate, and adverse gate jumps are inserted into
Corollary~\ref{cor:tube-step}. The input radius is $10^{-6}$, rounded upward for the scalar replay.
The certified radii at $3.7$ and $3.9$ are below
$1.199\cdot10^{-5}$ and $2.217\cdot10^{-5}$, respectively.
Every strict tube-invariance inequality must pass; the complete list is
\texttt{LOCAL\_FLOW\_ENCLOSURE.csv}, with its terminal report in
\texttt{LOCAL\_FLOW\_CERTIFICATE.json}. These computations certify that
the flow is within $3.5\cdot10^{-5}$ of the reference snapshots at $3.7$
and $3.9$. Section~\ref{sec:static} proves the static rate and geometry
bounds used here. The computational proof is conditional on the arithmetic
contract stated there; source code and exact binary inputs are supplied.

For completeness, the independent snapshot verification gives
\begin{align*}
 E(\xi_0,3.4)&\in[0.48314414,0.48315477],\\
 E(\xi_0,3.7)&\in[0.48298853,0.48299930],\\
 E(\xi_0,3.9)&\in[0.48306002,0.48307086],
 \qquad \xi_0=\xi_{-17\pi/36}.
\end{align*}
The displayed intervals are enlarged from the full-precision certified
values. They follow by Taylor's theorem for the full probe-error potential,
with error bound $\|g(y)\|R+H_gR^2/2$; all probe gates are constant
on the respective static balls. The first and second central gaps exceed
$0.00014484$ and $0.00006072$.

For any parameter vector in these balls,
$\sum_i\|w_i\|\|u_i\|\le\|\theta\|^2/2$.
Thus the full error on the unit circle is Lipschitz in the probe with
constant at most $(1+\|\theta\|^2/2)^2<9.084$.
Moving the probe by at most $10^{-6}$ in angle changes a difference of
two errors by less than $1.8168\cdot10^{-5}$. This proves the two
claimed uniform gaps. Continuity gives an interior minimum. The mass,
cone, and central derivative conclusions are the corresponding static
certificates. The terminal enclosure is smaller than $3\cdot10^{-5}$.
The speed bounds $0.385$ and $0.309$ ensure that extending either
endpoint forward by $10^{-5}$ stays in its radius-$3.5\cdot10^{-5}$
static ball, by a first-exit argument. Across the probe sector the masks
remain unchanged: the additional score movement is at most
$2.01\cdot10^{-6}$, below the unused central probe margin. On either ball
\[
 |\dot E(\xi)-\dot E(\xi_0)|
 \le2(1+P^2/2)P\|F\|\,\|\xi-\xi_0\|
 <5\cdot10^{-6},\qquad P<2.01.
\]
Subtract this cost from the static rate margins to obtain the asserted
uniform decreasing and increasing intervals. This also extends the
compatible flow through the short terminal interval.
\end{proof}

\paragraph{What the signed mechanism says.}
At both endpoints the strong-cluster contribution to the full OOD rate is
negative and the weak-cluster contribution is positive. The latter dominates
at $3.9$. Moreover its late lower bound exceeds its early upper bound, while
the strong-cluster contribution becomes less negative, by strict certified
endpoint margins. These are contributions to the \emph{same} gradient field
and full-network residual, not results of removing a cluster and retraining.
They explain the local reversal by signed competition between the two
training clusters. They do not prove that specialization alone is sufficient
or that it causes reversal for every specialized state.

\paragraph{The remaining initialized obligation.}
The certificate starts from the small neighborhood in
\eqref{eq:entry-missing}. It does not prove that the canonical initialized
flow reaches it. In particular, restarting a numerical or exact flow at
$\widehat\theta_{3.4}$ is not the same as certifying the path from the
archived seed. An initialized canonical theorem still requires
\eqref{eq:entry-missing}, or another proved entry enclosure for which the
local bounds close. The high-probability iid theorem requires an additional
probability argument and is not implied by this fixed-dataset result.
